% 夜半用户手册 · 第 5 章：录音与音频编辑（中文）
% T-DOC-001 逐章交付（CH05 任务未到前预写）。
% 事实来源：app.slint 菜单与回调、host.rs 录音接线、record_landing.rs、
% docs/manual/03-tracks-clips.md、cli.rs SHORTCUTS。

\chapter{录音与音频编辑}
\label{ch:recording-audio}

\classquote{感于物而动，故形于声。}{《礼记·乐记》}{录音是把当下的触动变成可编辑的声，后续剪辑与处理都为还原那份真切。}





这一章讲怎么把外部声音弄进工程：录音、导入、把音频片段摆到轨道上，以及
最基础的裁剪与淡化。录音与导入的素材都会进入工程资产池，随工程一起走。

\section{录音流程与输入监测}

\subsection{打开录音面板}

「文件 $\to$ 实时录音至 WAV...」打开录音面板。面板里有录音就绪（Arm）开关、
输入设备选择与输入测试；录音路径使用与工程一致的采样率（工程缺省 48 kHz）。
按键走带同样受录音状态影响：停录会结束写入，而不是继续偷偷录下去。

\screenshot{ch05-record-panel}{录音面板，显示就绪开关、输入设备选择与电平监测}

\subsection{录音落地}

停录成功后，录到的字节会按真实解码事实（帧数、采样率、声道、位深）落进
工程资产池，并在接线完整时向工程提交一次「加片段」操作——于是录音结果
直接出现在时间轴上，而不只是磁盘上的 WAV 文件。若当前装配没有把录音写入面
接进工程，应用只落盘、不假装提交过；这是无头或诊断形态的既有行为。

\begin{figure}[htbp]
\centering
\begin{tikzpicture}[node distance=1.2cm, >=stealth,
    box/.style={draw, rounded corners, align=center, minimum width=2.2cm, minimum height=0.9cm}]
  \node[box] (a) {选输入设备};
  \node[box, right=of a] (b) {Arm + 试听电平};
  \node[box, right=of b] (c) {录音 / 停录};
  \node[box, right=of c] (d) {字节进资产池};
  \node[box, right=of d] (e) {片段上时间轴};
  \draw[->] (a)--(b); \draw[->] (b)--(c); \draw[->] (c)--(d); \draw[->] (d)--(e);
\end{tikzpicture}
\caption{录音落地链路：选设备 $\to$ 就绪监测 $\to$ 录音/停录 $\to$ 资产池 $\to$ 时间轴片段}
\label{fig:ch05-record-flow}
\end{figure}

\section{音频片段的裁剪、淡入淡出与交叉淡化}

\subsection{导入与摆放}

「文件 $\to$ 导入音频文件...」把音频读进工程；「轨道 $\to$ 将最新导入音频
摆放到当前音频轨」把它放到当前音频轨的时间轴上。片段的落点由播放头与当前
选中轨决定，所以摆放前先确认播放头位置和选中轨，避免片段落到错误的轨或
位置。

\subsection{裁剪与淡化}

片段放好后，拖左右边缘可以改变它的时间范围：裁掉头部空白或尾部多余部分，
就是最常用的音频修整。淡入让片段开头从安静渐强，淡出让结尾渐弱；相邻两个
片段在重叠区间上做交叉淡化，可以让衔接不出现突变。

这些操作的具体入口（片段参数区、右键菜单或编辑菜单）以当前版本界面为准；
不同版本可能有差异，未出现的控件不要当作已提供。淡化只改变增益包络，不
改动资产池里的原始字节。

\screenshot{ch05-audio-clip-edges}{一个音频片段，标出左右可拖边缘与淡化手柄}

\subsection{不要被「完成」提示误导}

导出分轨时，量的是实际渲染与跳过情况；如果某条分轨里的音频片段全被跳过，
回执应如实给出片段跳过计数并将提示降级，而不是仍然只说「分轨完成」。看到
「完成」时，仍建议对照回执里的计数核对（详见第 8 章）。

\section{素材池与音频资产管理}

录音与导入得到的音频字节都进工程资产池，按内容哈希寻址。它的直接好处是
工程自带素材：拷走 \texttt{.yeban} 文件，用到的声音会随容器同行，换机器打开
不会变成无声。

\screenshot{ch05-asset-pool}{资产池面板，列出工程内音频资产与内容哈希}

\section{帮助主题映射}

\begin{table}[htbp]
\centering
\caption{本章小节 $\to$ 应用内帮助锚点}
\label{tab:ch05-help}
\begin{tabular}{ll}
\hline
本章小节 & 帮助主题 \\
\hline
打开录音面板 / 输入监测 & \texttt{record-panel-input} \\
录音落地 & \texttt{record-landing} \\
导入与摆放 & \texttt{audio-import-place} \\
裁剪与淡化 & \texttt{audio-trim-fade} \\
不要被「完成」提示误导 & \texttt{audio-receipt-honesty} \\
素材池与音频资产管理 & \texttt{audio-asset-pool} \\
\hline
\end{tabular}
\end{table}

\paragraph{本章学完能做什么} 完成一次录音并让片段落到目标轨上，导入音频、裁掉首尾、加淡入淡出，并通过回执核对导出结果。

\paragraph{动手练习} 新建一条音频轨，把播放头放到 001.01.000；用「实时录音至 WAV...」录两小节并停录，确认片段出现在时间轴上；裁掉开头空白，给结尾加淡出；再用「导入音频文件...」导入一个素材并用「将最新导入音频摆放到当前音频轨」放到第二条轨；最后导出分轨，对照回执里的片段计数。
